\newpage
\phantomsection
\label{prf:list-equality}

\begin{remark*}[Return]
\hyperref[def:list]{Return to Definition of List}
\end{remark*}

\begin{theorem*}[Equality of Lists]
Two lists $L = (x_1, \dots, x_n)$ and $M = (y_1, \dots, y_m)$ are equal if and only if $n = m$ and $x_j = y_j$ for all $j \in \{1, \dots, n\}$.
\end{theorem*}

\begin{proof}
\textbf{Professional Standard Proof.}~\\
A list of length $n$ is formally defined as a function $f: \{1, \dots, n\} \to S$. By the axiom of extensionality for functions, two functions $f$ and $g$ are equal if and only if they have the same domain and $f(j) = g(j)$ for all $j$ in that domain. Let $L$ be the function $j \mapsto x_j$ and $M$ be $j \mapsto y_j$. Then $L = M$ if and only if their domains $\{1, \dots, n\}$ and $\{1, \dots, m\}$ are equal (implying $n = m$) and $x_j = y_j$ for all $j \in \{1, \dots, n\}$.
\end{proof}

\begin{proof}
\textbf{Detailed Learning Proof.}~\\

\textbf{Step 1. Address the global constraint (Length).}
Before examining individual elements, we observe that the length $n$ is an intrinsic property of the list. If $n \neq m$, the lists are distinct objects in different spaces (e.g., $\mathbb{F}^n$ vs $\mathbb{F}^m$). Thus, $n=m$ is a necessary condition.

\textbf{Step 2. Address the local constraint (Positional Identity).}
Since lists are ordered, identity is not just about having the same "pool" of elements (like a set), but about the mapping of indices to values. We define equality element-wise:
\[ \forall j \in \{1, \dots, n\}, x_j = y_j \]
If there exists even one index $k$ such that $x_k \neq y_k$, the mapping is different, and thus the lists are not equal.

\textbf{Step 3. Critique of the Induction/Summation approach.}
While one could define a "difference sum" $D = \sum_{j=1}^n \delta(x_j, y_j)$ where $\delta$ returns $1$ if elements differ, this assumes the elements exist in a space where addition is defined. The functional definition used in Step 2 is more robust because it applies even if the list contains non-mathematical objects.
\end{proof}

\begin{remark*}[Proof structure]
The proof relies on the definition of a list as a finite sequence (function) and uses the principle of functional extensionality.
\end{remark*}

\begin{remark*}[Dependencies]~\\
\begin{itemize}
  \item \hyperref[def:list]{Definition of List}
  \item $\operatorname{FunctionEquality}$ (Volume I)
\end{itemize}
\end{remark*}
